\section{One Trie Descent for All Words}
\label{sec:trie}

A finite word can be shortened by deleting some of its letters. We need
a stronger conclusion than keeping its end state. Every comparison
created by the deletion has to be accounted for, including those with
other words of the tuple. \Cref{sec:refinement,sec:parsing,sec:contents}
prove a general shortening lemma for tuples of bounded depth
(\cref{lem:refinement}), by induction on the depth. \Cref{sec:buffers}
applies it to the stopped words and the stems of the rays.

\subsection{What a shortening preserves}
\label{sec:refinement}

We say a tuple \emph{refines} another tuple on a transition table if
it keeps each endpoint state and each of its cut states occurs in the
original coordinate. In addition, every comparison with nonempty target
rest must lift to such a comparison of the original words with the same
coordinates and pair of states. A comparison with empty target rest
lifts using the endpoint and cut-state conditions. Refinement therefore
preserves eligibility and safety.

To bound a refinement, we need to count the effects of a word read from
different states. Its \emph{word action} is the map sending each possible
entry state to the state after reading the word. When the input has
bounded depth, the following shortening depends on the number of these
maps.

\begin{lemma}[Short refinements]\label{lem:refinement}
Let a transition table have $k$ starts, $N$ states, and at most $K$
word actions. A tuple of depth at most $d$, with all coordinates
included in the depth test, has a refinement by coordinatewise
subsequences of total length at most
\[
 \FBound(k,N,K,d)
   = k\bigl(2(K+N+1)(2k+3)^{3d}+k+2\bigr)^{(d+1)^2}.
\]
\end{lemma}

The proof is an induction on the depth $d$. One \emph{level} of the
descent cuts all words of the tuple into blocks (\cref{sec:parsing}). The
lists of blocks are words over a new alphabet whose letters are blocks,
and they have smaller depth, so the induction shortens them. On the return,
the contents of the remaining blocks are shortened together
(\cref{sec:contents}). Three auxiliary readers carry the tests from one
level to the next. The \emph{transition reader} records at each cut the
transition that entered it, so that a test on blocks can refer to a cut
inside a block. The \emph{derived reader} reads a block list, one block per
letter, and has the same holes as the original reader on the expanded
words. The \emph{junction monitor} checks that a shortened block list still
expands to a word. \Cref{sec:buffers} adds a fourth one, the \emph{suffix
monitor}, which makes the shortening keep a prescribed end of each word.

\subsection{Parsing at markers}
\label{sec:parsing}

The last letter of any nonempty word is already one of the rests that
comparisons can see. We use the last letters of all words as
\emph{markers}, and cut every word after each occurrence of a marker.
A \emph{block} has a front marker, a content without markers, and a
closing marker. The first front marker of a word is a start symbol of its
coordinate. Every later front marker repeats the previous closing marker.
We call such lists \emph{linked}. Expanding a list writes each content
followed by its closing marker, so a repeated marker is written once.
\Cref{fig:trie} parses two words. Its letters belong to the alphabet of
some level, so at deeper levels they are blocks themselves.

\begin{figure}[t]
\centering
\begin{tikzpicture}[x=9mm,y=6mm,>=Stealth,
  txt/.style={font=\footnotesize}, nd/.style={circle,draw,fill=white,inner sep=1.5pt},
  mk/.style={fill=black!65,text=white,minimum width=5mm,inner sep=1.7pt,font=\footnotesize},
  ct/.style={draw=black!45,fill=black!9,minimum width=9mm,minimum height=4.8mm,
    inner sep=1.5pt,font=\footnotesize}]
  \node[txt,anchor=west] at (-.6,1.45) {Linked blocks};
  \node[txt] at (-.3,.5) {$w_1$};
  \node[txt] at (-.3,-.55) {$w_2$};
  \node[txt] at (.4,.5) {$\star_1$};
  \node[txt] at (.4,-.55) {$\star_2$};
  \node[ct] (cz) at (1.25,.5) {$cz$};
  \node[ct] (zo) at (1.25,-.55) {$z$};
  \node[mk] (l1) at (2.2,.5) {$\mu$};
  \node[mk] (l2) at (2.2,-.55) {$\mu$};
  \node[ct] (dz) at (3.15,.5) {$dz$};
  \node[ct] (em) at (3.15,-.55) {$\eps$};
  \node[mk] at (4.1,.5) {$\nu$};
  \node[mk] at (4.1,-.55) {$\mu$};
  \draw[->,shorten >=.5pt,shorten <=.5pt] (l1)--(dz);
  \draw[->,shorten >=.5pt,shorten <=.5pt] (l2)--(em);
  \node[txt,anchor=west] at (5.25,1.45) {Reverse trie of contents};
  \node[nd,label={[txt]right:$\eps$}] (root) at (9.9,0) {};
  \node[nd,label={[txt]below:$z$}] (z) at (7.8,0) {};
  \node[nd,label={[txt]left:$cz$}] (c) at (6,0.70) {};
  \node[nd,label={[txt]left:$dz$}] (d) at (6,-0.70) {};
  \draw[thick] (c)--node[txt,above]{$c$}(z);
  \draw[thick] (d)--node[txt,below]{$d$}(z);
  \draw[very thick] (z)--node[txt,above]{$z$}(root);
\end{tikzpicture}
\caption{The words $w_1=cz\mu dz\nu$ and $w_2=z\mu\mu$ end in $\nu$ and
$\mu$, so both letters are markers. The symbols $\star_1$ and $\star_2$ are
the start symbols of the two coordinates. A closing $\mu$ is also the front
marker of the next block. All contents ending in $z$ use the same edge to
the root.}
\label{fig:trie}
\end{figure}


Two nonempty rests of expanded words can form a comparison only if
the portions up to and including their next markers agree. The following
block lists must then be in the prefix relation. This reduces tests at
old cuts to tests on pairs of block occurrences.

The state after reading a block by itself would hide the states at
its interior cuts. We therefore use a transition reader. Its label at a
cut is the entering transition, the previous state and the letter just
read, and at the first cut it is the coordinate start. The derived reader
reads blocks as letters. It processes a block from its front marker
through its content. Its label at the cut after a block is the pair of the
state before the block and the block itself. From this pair one can
recompute the original label at every cut inside the block and at the cut
after its closing marker. The derived reader tests aligned cuts inside the
blocks and handles empty target rests separately. A last cut of the block
list may represent a proper cut of the original word, so both tests are
needed there.

A shortened list of blocks must still be linked. The junction monitor
remembers the previous closing marker and checks that it agrees with
the next front marker. It checks the first front marker against the
coordinate start. The product of the derived reader and this monitor
gives the reader for the next level.

Parsing lowers the depth. In a chain of two or more distinct nonempty
block rests, the first block is the same throughout the chain. Its
front marker is a marker of this level. A front marker that named a
coordinate start would force the rests to be the same whole list.
Prepending that front marker to each expanded rest recovers a chain
of rests of the original words. The marker by itself is another rest,
strictly shorter than all the others. At depth one this also implies
that each word parses into at most one block.

\subsection{Shortening the contents}
\label{sec:contents}

After the recursive step, we have fewer blocks to expand. They may still
have long contents. Deleting loops separately could make two different
contents end in the same sequence of retained letters, creating a new
comparison. We shorten the finite set of contents, the \emph{dictionary},
on one reverse trie. Its paths from the root spell the reversals of
contents.

Mark the root, the content endpoints, and the vertices where paths
branch. In each segment between marked vertices, read in the original
direction and delete loops of the word actions. The last letter of a
segment is retained. It distinguishes the outgoing branches in the
reverse trie. A shared segment is shortened once. Thus equal retained
suffixes have the same deleted gaps, and the state at each cut after a
retained letter is the same as at its original cut. This also keeps
the transitions into those cuts.

To finish the level, rename each surviving block by its shortened
content and keep its markers. The renaming is injective on the blocks
used. Every internal alignment lifts to an old alignment, so the
renamed block lists still pass all the tests. They expand to a valid
tuple at this level.

To obtain a refinement of a given tuple $W$, we run the descent on a
transition reader built from $W$. It accepts exactly the endpoint labels
of $W$. For each ordered pair of coordinates, it admits a pair of labels at
a comparison with nonempty target rest if no such comparison of $W$
between the same coordinates carries them. It admits a host label against
an endpoint target if the label does not occur in the host coordinate of
$W$. Then $W$ is a hole of this reader. In every hole of it, the endpoint
labels, the cut labels, and the comparisons with nonempty target rest
occur in the same coordinates of $W$. The short hole of this reader is
therefore the required refinement. \Cref{app:trie} proves the transfer at
each level and counts the support and length on expansion.

\subsection{Keeping the buffers}
\label{sec:buffers}

Cutting off each ray at its stem would discard comparisons using its
periodic part. The \emph{onset} of a ray is the least cut from which it
is periodic, and its \emph{stem} is the prefix before the onset. We
attach a finite part of the periodic text to the stem, its \emph{buffer}
of the next $B=(C+N+3)D+4$ letters. We show below that the buffer is long
enough to contain every comparison with a stopped target. A new terminal
marker follows the buffer. It has the identity action in the reader.

The tuple for shortening consists of these marked words and the stopped
words. Its depth is at most $C+1$. Any nonempty rest of a marked word
ends at its marker, so among the prefixes of a test word these rests
contribute only the length up to the first terminal marker. The stopped
words contribute at most $C$ lengths.

We keep the \emph{terminal window} of each word, its last $B+2$ letters
or the whole word if shorter, and the states at the cuts of this window.
This keeps the stem's last letter when it exists, the buffer, and the
terminal marker. It also keeps the state reached at each onset. To impose
this requirement we run the reader with a suffix monitor. Its states
record the longest current suffix that is a factor of a terminal window,
together with the reader state before that suffix, in the style of the
automaton of Aho and Corasick \cite{AhoCorasick75}. A step updates the
record. The final record gives the prescribed suffixes and their states.
The proof in \cref{app:trie} bounds the word actions of this monitor
and applies one refinement to the marked tuple.

\begin{figure}[t]
\centering
\begin{tikzpicture}[x=9mm,y=6.7mm,>=Stealth,
  leg/.style={font=\footnotesize}, seg/.style={draw,fill=black!10},
  buff/.style={draw,fill=black!28}, kept/.style={draw,fill=black!65},
  bmark/.style={fill=black!85,text=white,font=\footnotesize,inner sep=1.4pt},
  brace/.style={decorate,decoration={brace,amplitude=3pt}}]
  \node[leg,anchor=east] at (-.25,0) {old ray};
  \node[leg,anchor=east] at (1.35,-1.55) {shortened ray};
  \draw[seg] (0,-.18) rectangle (3.65,.18);
  \draw[kept] (3.65,-.18) rectangle (4.1,.18);
  \draw[buff] (4.1,-.18) rectangle (7.25,.18);
  \draw[seg] (1.6,-1.73) rectangle (3.65,-1.37);
  \draw[kept] (3.65,-1.73) rectangle (4.1,-1.37);
  \draw[buff] (4.1,-1.73) rectangle (7.25,-1.37);
  \foreach \y in {0,-1.55}{
    \draw[->] (7.25,\y)--(10.7,\y);
    \foreach \x in {5.15,6.2,7.25,8.30,9.35,10.4}
      \draw[black!70] (\x,\y-.18)--(\x,\y+.18);
    \draw[densely dotted] (7.25,\y+.18)--(7.25,\y+.37);
    \node[bmark,anchor=south] at (7.25,\y+.37) {$\#$};
  }
  \draw[brace] (0,.32)--node[leg,above=2pt]{stem}(4.1,.32);
  \draw[brace] (4.1,.32)--node[leg,above=2pt]{buffer}(7.25,.32);
  \node[leg,above] at (4.1,-1.31) {onset};
  \draw[very thick] (9.85,-1.89)--(9.85,-1.18);
  \node[leg,below] at (9.85,-1.92) {truncation};
\end{tikzpicture}
\caption{The shortening sees finite prefixes ending at the marker $\#$.
It keeps the last stem letter (dark) and the buffer. The rays continue
with the original periodic text when truncation lengths are chosen.}
\label{fig:ray}
\end{figure}


Two properties transfer to the new rays. First, comparisons with a
stopped target fit inside the buffered prefix. After the onset, the state
at a cut and its \emph{phase}, the distance from the onset modulo the
period, determine all later states and letters. So the host cut can be
moved back to the first cut with the same state and phase. A match from
there that extended beyond the buffer would give a stopped word a periodic
suffix that is too long for its depth. Second, matches through the new
stems cannot create admitted equal infinite rests that violate the
offsets. If equal infinite rests start before the onsets, they have the
same distance to their onsets. Their finite rests through the terminal
markers lift together. If the cuts are in the periodic parts, the states
and rests come from the old rays directly. The onsets themselves are
kept. Since the onset is the least cut from which the ray is periodic,
the last stem letter differs from the letter one period later. Keeping it
makes the shortened ray become periodic exactly at its onset.

\begin{lemma}[Shortening a periodic tuple]\label{lem:buffered}
For a periodic tuple whose stopped words have depth at most $C$ and whose
periods are at most $D$ from \eqref{eq:period-bound}, let
\begin{align*}
  B&=(C+N+3)D+4, & M&=1+k(B+2)^2,\\
  S&=\FBound\bigl(k,MN,M(K+1),C+1\bigr).
\end{align*}
One trie descent gives a periodic tuple with each stopped word and each
stem of length at most $S$, and with periods at most $D$.
Here $K$ bounds the word actions of the original reader.
\end{lemma}

The buffer safety and offset proofs are in \cref{app:truncation}.
The finite parts of the periodic tuple are now bounded.
We still have to find finite truncations of the new rays.
